\section{عدم الانعدام عند الواحد للشخصيات غير الحقيقية}

نغلق الآن الحالة الأولى من بوابة عدم الانعدام. تعتمد الحجة على متراجحة
مثلثية أولية، لكنها تستعملها داخل لوغاريتم حاصل ضرب أويلري مضبوط؛ فلا
يجوز اختزالها إلى عبارة غير مبرهنة من نوع «المنتج الأويلري موجب».

\begin{lemma}[متراجحة ديريشليه المثلثية]
\label{lem:dirichlet-trigonometric-inequality}
\resultid{ANT-LEM-07-01}
\provedhere
لكل \(\theta\in\mathbb R\):
\[
3+4\cos\theta+\cos(2\theta)
=2(1+\cos\theta)^2\geq0.
\]
\end{lemma}

\begin{proof}
باستعمال \(\cos(2\theta)=2\cos^2\theta-1\) نحصل على
\[
3+4\cos\theta+\cos(2\theta)
=2+4\cos\theta+2\cos^2\theta
=2(1+\cos\theta)^2.
\]
\end{proof}

\begin{theorem}[عدم الانعدام للشخصية غير الحقيقية]
\label{thm:nonreal-character-nonvanishing-at-one}
\resultid{ANT-THM-07-10}
\provedhere
إذا كانت \(\chi\) شخصية ديريشليه غير حقيقية وغير رئيسية، فإن
\[
L(1,\chi)\ne0.
\]
\end{theorem}

\begin{proof}
لـ\(\sigma>1\) نعرف الكمية الحقيقية الموجبة
\[
F(\sigma)
=
\zeta(\sigma)^3
|L(\sigma,\chi)|^4
|L(\sigma,\chi^2)|.
\]
جميع العوامل غير منعدمة في \(\sigma>1\)، ولذلك يمكن أخذ اللوغاريتمات
الحقيقية للقيم المطلقة. ومن التوسعات الأويلرية المطلقة نحصل على
\begin{align*}
\log F(\sigma)
&=
\sum_p\sum_{k=1}^{\infty}
\frac{3+4\Re(\chi(p)^k)+\Re(\chi(p)^{2k})}
{k p^{k\sigma}}.
\end{align*}
إذا كانت \(\chi(p)=0\)، فالبسط يساوي \(3\). وإذا كانت
\(|\chi(p)|=1\)، فاكتب \(\chi(p)^k=e^{i\theta}\). عندئذ يصبح البسط
\[
3+4\cos\theta+\cos(2\theta)\geq0
\]
بحسب اللمّة السابقة. إذن
\[
\log F(\sigma)\geq0,
\qquad
F(\sigma)\geq1
\quad(\sigma>1).
\]

نفترض خلاف المطلوب أن \(L(1,\chi)=0\)، ولتكن رتبة الصفر \(m\geq1\).
بما أن
\[
L(\overline s,\overline\chi)
=
\overline{L(s,\chi)},
\]
فإن القيمة على المحور الحقيقي تحقق
\(|L(\sigma,\chi)|\asymp(\sigma-1)^m\) عندما
\(\sigma\to1^+\).

ولأن \(\chi\) غير حقيقية، فإن \(\chi^2\) غير رئيسية: فلو كانت
\(\chi^2=\chi_0\) على الوحدات لكانت كل قيمة لـ\(\chi\) مساوية
لـ\(\pm1\)، فتكون \(\chi\) حقيقية. ومن ثم \(L(s,\chi^2)\) دالة تامة،
وبخاصة محدودة قرب \(s=1\). كذلك
\[
\zeta(\sigma)\asymp\frac1{\sigma-1}.
\]
وعليه
\[
F(\sigma)
\ll
(\sigma-1)^{-3}(\sigma-1)^{4m}
=
(\sigma-1)^{4m-3}
\longrightarrow0
\]
لأن \(m\geq1\). وهذا يناقض \(F(\sigma)\geq1\). إذن
\(L(1,\chi)\ne0\).
\end{proof}

\begin{remark}[موضع استعمال كون الشخصية غير حقيقية]
الخطوة الحاسمة هي أن \(\chi^2\) غير رئيسية، ومن ثم لا يضيف
\(L(s,\chi^2)\) قطبًا عند \(s=1\). إذا كانت \(\chi\) حقيقية غير
رئيسية، فإن \(\chi^2=\chi_0\) على الوحدات، فتتغير موازنة رتب الأقطاب
والأصفار، ولا تنتج الحجة السابقة تناقضًا. لذلك تبقى الحالة الحقيقية
بوابة مستقلة.
\end{remark}

\begin{corollary}
إذا كانت \(\chi\) غير حقيقية، فإن الصفر البديهي عند \(s=0\) في الحالة
الزوجية بسيط، لأن المعادلة الوظيفية تربطه بقيمة
\(L(1,\overline\chi)\ne0\).
\end{corollary}

\section{حالة بوابة عدم الانعدام بعد إغلاق الحالة المركبة}

أصبحت الحالة غير الحقيقية مثبتة داخل الموسوعة. أما النتيجة العامة
\[
L(1,\chi)\ne0
\]
لكل شخصية غير رئيسية فتبقى بحالة \textenglish{DEFERRED} إلى أن يكتمل
برهان الشخصيات الحقيقية من دون استعمال دائري لمبرهنة ديريشليه للأوليات
في المتتاليات الحسابية.
